%【紙面設定】
%b4,横書き,2段
\documentclass[paper=b4j,landscape,twocolumn,fleqn]{jlreq}
%横書き用の調整
\special{papersize=\the\paperwidth,\the\paperheight}
%間に線をいれる
\setlength{\columnseprule}{0.4pt}
%ページ番号消す
\pagestyle{empty}
%ページ余白の調整
\usepackage[margin=10truemm]{geometry}
%画像挿入用のパッケージ
\usepackage[dvipdfmx]{graphicx}
%数学用のパッケージ
\usepackage{amsmath}
% 見出しの設定：行取りを1行にし、前後のスペースを極限まで削る
\ModifyHeading{section}{
  lines=1,
  before_space=10pt, % セクションの上の余白（適度に残すと見やすい）
  after_space=0pt    % セクションのしたの余白をゼロにする
}


%【本文】
\begin{document}
%1段落目
\section{五則演算(÷は/で表す.×は*で表す.)}
\begin{align*}
    5+8 &= & 9-4 &= & 2*3 &= & 6/2 &= & 5\%2 &= \\
    -9+4 &= & 3-10 &= & -4*5 &= & -10/5 &= & 10\%3 &= \\
    15+27 &= & 50-18 &= & 12*4 &= & 48/4 &= & 25\%4 &= \\
    -50+12 &= & 20-85 &= & 15*-3 &= & 100/-25 &= & 100\%7 &= \\
    1.2+3.8 &= & 4.5-1.2 &= & 0.5*6 &= & 4.8/2 &= & 50\%6 &= \\
    -4.5+1.2 &= & 0.8-3.5 &= & -1.2*4 &= & -7.2/0.9 &= & 80\%9 &= \\
    125+380 &= & 500-125 &= & 25*12 &= & 144/12 &= & 123\%10 &= \\
    -600+250 &= & 100-456 &= & 14*-10 &= & -117/9 &= & 250\%15 &= \\
    0.45+0.55 &= & 1.25-0.75 &= & 0.8*0.5 &= & 0.75/0.25 &= & 1000\%3 &= \\
    -12.5+7.5 &= & 1000-1500 &= & -15*1.2 &= & 102/-3 &= & 1024\%10 &= 
\end{align*}
\section{ルートと乗数の計算}
\begin{align*}
    2^2 &= & 2^2 * 2^1 &= & \sqrt{4} &= & \sqrt{4} * \sqrt{9} &= \\
    3^2 &= & 3^3 / 3^1 &= & \sqrt{9} &= & \sqrt{2} * \sqrt{8} &= \\
    2^4 &= & 2^2 * 2^2 &= & \sqrt{16} &= & \sqrt{12} / \sqrt{3} &= \\
    5^2 &= & 10^2 / 10^1 &= & \sqrt{25} &= & \sqrt{3} * \sqrt{12} &= \\
    4^3 &= & 2^3 * 2^2 &= & \sqrt{49} &= & \sqrt{18} / \sqrt{2} &= \\
    6^2 &= & 5^3 / 5^2 &= & \sqrt{64} &= & \sqrt{2} * \sqrt{32} &= \\
    10^2 &= & 2^2 * 2^4 &= & \sqrt{100} &= & \sqrt{50} / \sqrt{2} &= \\
    12^2 &= & 3^4 / 3^2 &= & \sqrt{144} &= & \sqrt{72} / \sqrt{2} &= 
\end{align*}
\section{分数の計算}
\begin{align*}
    \frac{1}{4} + \frac{1}{4} &= & \frac{4}{5} - \frac{1}{5} &= & \frac{1}{3} \times \frac{1}{2} &= & \frac{2}{5} \div 2 &= \\
    -\frac{2}{3} + \frac{1}{3} &= & \frac{1}{4} - \frac{3}{4} &= & -\frac{1}{2} \times \frac{1}{4} &= & \frac{1}{5} \div (-2) &= \\
    \frac{1}{2} + \frac{1}{3} &= & \frac{1}{2} - \frac{1}{4} &= & \frac{2}{3} \times \frac{3}{5} &= & \frac{3}{4} \div 2 &= \\
    -\frac{1}{2} + \frac{1}{6} &= & \frac{1}{6} - \frac{1}{2} &= & \frac{2}{5} \times \left(-\frac{1}{4}\right) &= & -\frac{3}{4} \div 3 &= \\
    \frac{3}{4} + \frac{1}{6} &= & \frac{5}{6} - \frac{1}{4} &= & \frac{5}{8} \times \frac{4}{15} &= & \frac{2}{3} \div \frac{1}{2} &= \\
    -\frac{1}{2} - \frac{1}{3} - \frac{1}{6} &= & \frac{1}{10} - \frac{4}{15} &= & -\frac{15}{16} \times \frac{4}{5} &= & 10 \div \left(-\frac{5}{2}\right) &= 
\end{align*}

%2段落目
\newpage
\section{五則演算(÷は/で表す.×は*で表す.)}
\begin{align*}
    5+8 &= & 9-4 &= & 2*3 &= & 6/2 &= & 5\%2 &= \\
    -9+4 &= & 3-10 &= & -4*5 &= & -10/5 &= & 10\%3 &= \\
    15+27 &= & 50-18 &= & 12*4 &= & 48/4 &= & 25\%4 &= \\
    -50+12 &= & 20-85 &= & 15*-3 &= & 100/-25 &= & 100\%7 &= \\
    1.2+3.8 &= & 4.5-1.2 &= & 0.5*6 &= & 4.8/2 &= & 50\%6 &= \\
    -4.5+1.2 &= & 0.8-3.5 &= & -1.2*4 &= & -7.2/0.9 &= & 80\%9 &= \\
    125+380 &= & 500-125 &= & 25*12 &= & 144/12 &= & 123\%10 &= \\
    -600+250 &= & 100-456 &= & 14*-10 &= & -117/9 &= & 250\%15 &= \\
    0.45+0.55 &= & 1.25-0.75 &= & 0.8*0.5 &= & 0.75/0.25 &= & 1000\%3 &= \\
    -12.5+7.5 &= & 1000-1500 &= & -15*1.2 &= & 102/-3 &= & 1024\%10 &= 
\end{align*}
\section{ルートと乗数の計算}
\begin{align*}
    2^2 &= & 2^2 * 2^1 &= & \sqrt{4} &= & \sqrt{4} * \sqrt{9} &= \\
    3^2 &= & 3^3 / 3^1 &= & \sqrt{9} &= & \sqrt{2} * \sqrt{8} &= \\
    2^4 &= & 2^2 * 2^2 &= & \sqrt{16} &= & \sqrt{12} / \sqrt{3} &= \\
    5^2 &= & 10^2 / 10^1 &= & \sqrt{25} &= & \sqrt{3} * \sqrt{12} &= \\
    4^3 &= & 2^3 * 2^2 &= & \sqrt{49} &= & \sqrt{18} / \sqrt{2} &= \\
    6^2 &= & 5^3 / 5^2 &= & \sqrt{64} &= & \sqrt{2} * \sqrt{32} &= \\
    10^2 &= & 2^2 * 2^4 &= & \sqrt{100} &= & \sqrt{50} / \sqrt{2} &= \\
    12^2 &= & 3^4 / 3^2 &= & \sqrt{144} &= & \sqrt{72} / \sqrt{2} &= 
\end{align*}
\section{分数の計算}
\begin{align*}
    \frac{1}{4} + \frac{1}{4} &= & \frac{4}{5} - \frac{1}{5} &= & \frac{1}{3} \times \frac{1}{2} &= & \frac{2}{5} \div 2 &= \\
    -\frac{2}{3} + \frac{1}{3} &= & \frac{1}{4} - \frac{3}{4} &= & -\frac{1}{2} \times \frac{1}{4} &= & \frac{1}{5} \div (-2) &= \\
    \frac{1}{2} + \frac{1}{3} &= & \frac{1}{2} - \frac{1}{4} &= & \frac{2}{3} \times \frac{3}{5} &= & \frac{3}{4} \div 2 &= \\
    -\frac{1}{2} + \frac{1}{6} &= & \frac{1}{6} - \frac{1}{2} &= & \frac{2}{5} \times \left(-\frac{1}{4}\right) &= & -\frac{3}{4} \div 3 &= \\
    \frac{3}{4} + \frac{1}{6} &= & \frac{5}{6} - \frac{1}{4} &= & \frac{5}{8} \times \frac{4}{15} &= & \frac{2}{3} \div \frac{1}{2} &= \\
    -\frac{1}{2} - \frac{1}{3} - \frac{1}{6} &= & \frac{1}{10} - \frac{4}{15} &= & -\frac{15}{16} \times \frac{4}{5} &= & 10 \div \left(-\frac{5}{2}\right) &= 
\end{align*}
\end{document}


